%% privacy-policy.tex
%% Template per Informativa Privacy — PoliNetwork APS
%%
%% Compila con: xelatex privacy-policy.tex
%% (dalla root del progetto)
%%
\documentclass[legal]{../core/polinetwork}

%% ====================================================================
%% DATI DOCUMENTO — Modifica qui sotto
%% ====================================================================

\pnTitle{Informativa sul Trattamento dei Dati Personali}
\pnSubtitle{ai sensi del Regolamento UE 2016/679 (GDPR)}
\pnDate{1 aprile 2026}
\pnVersion{1.0}

% Soggetto specifico dell'informativa (modifica in base al contesto)
% Es: "per i candidati al Recruitment", "per gli utenti del sito web", ecc.
\pnSubject{Informativa privacy per [specificare contesto]}

%% ====================================================================

\begin{document}
\thispagestyle{firstpage}

%% Header istituzionale
\pnMakeHeader

%% Titolo e metadati
\pnMakeMetadata

%% ====================================================================
%% CONTENUTO — Incolla/modifica il testo qui sotto
%% ====================================================================

La presente informativa descrive le modalità di trattamento dei dati personali
raccolti da \textbf{PoliNetwork APS} in relazione a [specificare contesto].

%% --- Sezione 1 ---
\pnNumberedSection{Titolare del trattamento}

Ai sensi degli artt.\,4 e 24 del Reg.\,UE 2016/679, il titolare del trattamento
è individuabile nella figura del legale rappresentante dell'associazione
\textbf{PoliNetwork APS}, il suo Presidente.

L'associazione non ha nominato un Responsabile della Protezione dei Dati (DPO)
in quanto non ricorrono le condizioni di cui all'art.\,37 del GDPR.

\pnSubsection{Contatti per esercitare i diritti privacy}

\begin{itemize}
  \item Indirizzo email: \href{mailto:privacy@polinetwork.org}{privacy@polinetwork.org}
        (in alternativa \href{mailto:direttivo@polinetwork.org}{direttivo@polinetwork.org})
  \item PEC: \href{mailto:polinetwork@pec.it}{polinetwork@pec.it}
\end{itemize}

%% --- Sezione 2 ---
\pnNumberedSection{Tipologia di dati raccolti}

% Inserisci qui l'elenco dei dati raccolti, ad esempio:
Tramite [specificare strumento/servizio] raccogliamo e trattiamo i seguenti dati personali:

\begin{itemize}
  \item \textbf{Dati anagrafici e di contatto}: nome, cognome, indirizzo email, ecc.
  \item \textbf{Dati tecnici}: indirizzo IP, browser, sistema operativo, ecc.
  \item \textbf{Dati di utilizzo}: pagine visitate, tempo di permanenza, ecc.
\end{itemize}

%% --- Sezione 3 ---
\pnNumberedSection{Finalità e base giuridica del trattamento}

I dati personali saranno trattati esclusivamente per le seguenti finalità:

\begin{enumerate}
  \item [Finalità 1];
  \item [Finalità 2];
  \item [Finalità 3].
\end{enumerate}

La base giuridica del trattamento è [specificare: consenso esplicito / legittimo
interesse / esecuzione contrattuale / obbligo legale] ai sensi dell'art.\,6,
par.\,1, lett.\,[a/b/c/f] del GDPR.

%% --- Sezione 4 ---
\pnNumberedSection{Modalità del trattamento}

Il trattamento dei dati personali avviene mediante strumenti manuali, informatici
e telematici. Le logiche applicate sono strettamente correlate alle finalità stesse
e implementate in modo da garantire la massima sicurezza e riservatezza dei dati,
prevenendone la perdita, gli usi illeciti e gli accessi non autorizzati.

%% --- Sezione 5 ---
\pnNumberedSection{Conservazione dei dati}

I dati personali saranno conservati per il tempo strettamente necessario a
conseguire le finalità per cui sono stati raccolti. Nello specifico:

\begin{itemize}
  \item [Specificare durata e condizioni di conservazione];
  \item al termine del periodo indicato, i dati saranno cancellati o resi
        completamente anonimi.
\end{itemize}

%% --- Sezione 6 ---
\pnNumberedSection{Destinatari e accesso ai dati}

I dati personali non saranno comunicati a soggetti terzi esterni all'associazione
né diffusi pubblicamente, salvo quanto di seguito specificato.

\pnSubsection{Trasferimento verso Paesi terzi}

% Inserisci qui le informazioni sui trasferimenti extra-UE se applicabili
Qualora i dati vengano trasferiti verso Paesi terzi, ciò avverrà nel rispetto
del GDPR sulla base di [specificare garanzie: decisione di adeguatezza,
clausole contrattuali standard, ecc.].

%% --- Sezione 7 ---
\pnNumberedSection{Diritti degli interessati}

L'interessato potrà far valere i propri diritti come espressi dagli artt.\,15--22
del Regolamento UE 2016/679, rivolgendosi al Titolare del Trattamento tramite
i recapiti indicati nella Sezione 1 del presente documento.

Nello specifico, l'interessato ha il diritto, in qualunque momento, di:

\begin{itemize}
  \item \textbf{Diritto di accesso} (art.\,15): ottenere la conferma
    dell'esistenza o meno di dati personali che lo riguardano;
  \item \textbf{Diritto di rettifica} (art.\,16): richiedere l'aggiornamento
    o la rettifica di dati inesatti;
  \item \textbf{Diritto alla cancellazione} (art.\,17): ottenere la cancellazione
    dei propri dati;
  \item \textbf{Diritto di limitazione} (art.\,18): richiedere che i dati siano
    unicamente conservati in determinate ipotesi;
  \item \textbf{Diritto alla portabilità} (art.\,20): ricevere i dati in formato
    strutturato e leggibile da dispositivo automatico;
  \item \textbf{Diritto di opposizione} (art.\,21): opporsi al trattamento
    per motivi connessi alla propria situazione particolare.
\end{itemize}

\pnHighlightBox{%
  Fatto salvo ogni altro ricorso amministrativo e giurisdizionale,
  l'interessato ha il diritto di proporre reclamo al Garante per la
  protezione dei dati personali ai sensi dell'art.\,77 del GDPR
  (\href{https://www.garanteprivacy.it}{www.garanteprivacy.it}).
}

%% ====================================================================
%% FOOTER CONTATTI
%% ====================================================================

\vfill
\pnContactBlock

\end{document}
